%--------------------------------------------------------------------------------------------------------------------------
%                    The range interface
%--------------------------------------------------------------------------------------------------------------------------
\section{The lazy range interface: iRange}
\index{iRange}
\index{range}
\renewcommand{\container}{Range}

A range is a read-only, lazy recipe for visiting values from a container or a
C array.  Adaptors such as \verb|Filter|, \verb|Transform| and \verb|Take|
do not copy or traverse the source.  Work starts only when a cursor is opened
or a terminal algorithm is called.  This is similar to a C++ ranges pipeline,
but the C interface makes status, ownership and element sizes explicit.

The source storage is borrowed.  It must outlive the range and every cursor
opened from it.  Finalizing a range destroys only the range recipe; it never
destroys the source container or array.  A pipeline may be evaluated more than
once because each terminal operation opens an independent cursor.

Include \verb|range.h| directly, or include \verb|containers.h|, which also
exposes this interface.

\subsection{Callback and interface types}

\begin{verbatim}
typedef struct Range Range;
typedef struct RangeCursor RangeCursor;

typedef int (*RangePredicate)(const void *element, void *arg);
typedef int (*RangeTransformFunction)(const void *input, void *output,
                                      void *arg);
typedef int (*RangeFoldFunction)(void *accumulator, const void *element,
                                 void *arg);
typedef int (*RangeVisitFunction)(const void *element, void *arg);

typedef struct tagRangeInterface {
    int (*FromSequential)(SequentialContainer *source, Range **result);
    int (*FromSequentialWithAllocator)(
        SequentialContainer *source,
        const ContainerAllocator *allocator, Range **result);
    int (*FromGeneric)(GenericContainer *source, size_t elementSize,
                       Range **result);
    int (*FromGenericWithAllocator)(
        GenericContainer *source, size_t elementSize,
        const ContainerAllocator *allocator, Range **result);
    int (*FromArray)(const void *data, size_t count, size_t elementSize,
                     Range **result);
    int (*FromArrayWithAllocator)(const void *data, size_t count,
                                  size_t elementSize,
                                  const ContainerAllocator *allocator,
                                  Range **result);

    int (*Filter)(Range **range, RangePredicate predicate, void *arg);
    int (*Transform)(Range **range, size_t outputElementSize,
                     RangeTransformFunction transform, void *arg);
    int (*Take)(Range **range, size_t count);
    int (*Drop)(Range **range, size_t count);
    int (*TakeWhile)(Range **range, RangePredicate predicate,
                     void *arg);
    int (*DropWhile)(Range **range, RangePredicate predicate,
                     void *arg);
    int (*Concat)(Range **left, Range **right);

    int (*Open)(const Range *range, RangeCursor **result);
    int (*Next)(RangeCursor *cursor, const void **element);
    int (*DeleteCursor)(RangeCursor *cursor);

    size_t (*GetElementSize)(const Range *range);
    ErrorFunction (*SetErrorFunction)(Range *range,
                                      ErrorFunction function);

    int (*ForEach)(const Range *range, RangeVisitFunction function,
                   void *arg);
    int (*Fold)(const Range *range, void *accumulator,
                RangeFoldFunction function, void *arg);
    int (*FindIf)(const Range *range, RangePredicate predicate,
                  void *arg, void *result);
    int (*AnyOf)(const Range *range, RangePredicate predicate,
                 void *arg);
    int (*AllOf)(const Range *range, RangePredicate predicate,
                 void *arg);
    int (*CountIf)(const Range *range, RangePredicate predicate,
                   void *arg, size_t *result);
    int (*ToVector)(const Range *range, Vector **result);

    int (*Finalize)(Range *range);
} RangeInterface;

extern RangeInterface iRange;
\end{verbatim}

\subsection{Return and callback protocol}

Range functions use signed integer status values so that normal termination
cannot be confused with an error:

\begin{center}
\begin{tabular}{|c|p{10.5cm}|}
\hline
\textbf{Status} & \textbf{Meaning} \\\hline
positive & Success, an available element, a true predicate, or continuation. \\\hline
zero & Normal end, not found, false, or a visitor request to stop. \\\hline
negative & A \verb|CONTAINER_ERROR_*| value. \\\hline
\end{tabular}
\end{center}

A \verb|RangePredicate| returns positive for true, zero for false, or a
negative error.  A \verb|RangeVisitFunction| returns positive to continue,
zero to stop normally, or a negative error.  Transform and fold callbacks
must return a positive value on success.  Returning zero from either is a
protocol violation, represented by:
\begin{center}
\verb|CONTAINER_ERROR_WRONGELEMENT|
\end{center}
Negative callback errors pass through unchanged.

The interface reports errors through both the return value and an error
function.  New ranges initially use \verb|iError.RaiseError|.
\verb|SetErrorFunction| installs a handler on the complete pipeline and
returns the former handler.  Invalid calls that have no usable range instance
are reported through the global \verb|iError| handler.  Applications should
always check negative return values even when they replace the reporting
callback.

Errors encountered while advancing a cursor are sticky.  Once \verb|Next|
returns a negative value, later calls on that cursor return the same value.
The cursor must still be released with \verb|DeleteCursor|.

\subsection{Sources}

\verb|FromSequential| reads the element size from the sequential interface.
The generic adapter \verb|iSequentialContainer.GetElementSize| is supported
by List, Dlist and Vector.  A sequential implementation that cannot expose a
fixed element size is rejected with \verb|CONTAINER_ERROR_INCOMPATIBLE|.

\verb|FromGeneric| can adapt another iterator-based container when the caller
supplies its fixed \verb|elementSize|.  An element size of zero denotes an
opaque element representation.  Opaque ranges can be traversed and tested,
but operations that must copy an element, including \verb|FindIf| and
\verb|ToVector|, return \verb|CONTAINER_ERROR_INCOMPATIBLE|.

\verb|FromArray| reads \verb|count| adjacent elements of \verb|elementSize|
bytes starting at \verb|data|.  A null data pointer is accepted only when the
count is zero.  The array is borrowed, not copied.

The \verb|WithAllocator| variants use the supplied allocator for all recipe,
cursor and transformation storage derived from that source.  The allocator
object and its functions must remain valid until the range and its cursors are
destroyed.

All source constructors set \verb|*result| to \Null on failure.  They return
a positive value on success and a negative error code otherwise.

\subsection{Adaptors and ownership}

The unary adaptors build the following lazy operations:

\begin{center}
\begin{tabular}{|l|p{10.2cm}|}
\hline
\verb|Filter| & Emits elements for which the predicate is true. \\\hline
\verb|Transform| & Produces one output value of
\verb|outputElementSize| bytes for each input value. \\\hline
\verb|Take| & Emits at most the first \verb|count| values. \\\hline
\verb|Drop| & Skips at most the first \verb|count| values. \\\hline
\verb|TakeWhile| & Emits values until its predicate first becomes false. \\\hline
\verb|DropWhile| & Skips values while its predicate is true, then emits all
remaining values. \\\hline
\end{tabular}
\end{center}

Each unary adaptor accepts \verb|Range **|.  On success it replaces the
handle with the new outer recipe node, which owns the previous node.  On any
failure the original handle and pipeline remain unchanged.  Range handles
therefore have unique ownership: do not copy a handle and later adapt, open,
or finalize the retained alias.  Finalize only the current outermost handle.

\verb|Concat| requires two ranges with the same element size.  On success it
replaces \verb|*left| with the combined pipeline and sets \verb|*right| to
\Null.  On failure both handles remain unchanged.  This makes allocation and
element-size failures transactional and gives the caller an unambiguous
cleanup path.

\subsection{Cursors}

\verb|Open| creates a new cursor for a pipeline.  \verb|Next| returns one and
places an element pointer in \verb|*element|, returns zero and sets the pointer
to \Null at normal end, or returns a negative error.  The pointer is borrowed
from the cursor and remains valid only until the next call to \verb|Next| or
until \verb|DeleteCursor|.  Copy the value when it must be retained.

Container-backed cursors detect structural source changes.  If the container
changes after the cursor is opened, \verb|Next| returns
\verb|CONTAINER_ERROR_OBJECT_CHANGED| and the error remains sticky.  Do not
modify source storage while a cursor or terminal operation is using it.

For example, manual cursor traversal keeps end and failure distinct:

\begin{verbatim}
RangeCursor *cursor = NULL;
const void *element;
int status = iRange.Open(range, &cursor);

while (status > 0 && (status = iRange.Next(cursor, &element)) > 0) {
    const int value = *(const int *)element;
    /* use value */
}

if (cursor != NULL) {
    int deleteStatus = iRange.DeleteCursor(cursor);
    if (status >= 0)
        status = deleteStatus;
}
/* status == 0 means normal end; status < 0 is an error. */
\end{verbatim}

\subsection{Terminal algorithms}

Terminal algorithms evaluate a pipeline using their own independent cursor
and do not consume or modify the range recipe.

\begin{center}
\begin{tabular}{|l|p{10.2cm}|}
\hline
\verb|ForEach| & Visits values until the source ends or the visitor returns
zero.  Source exhaustion returns one; a visitor-requested stop returns zero.
Both are normal outcomes. \\\hline
\verb|Fold| & Calls the fold function for every value, updating the
caller-owned accumulator. \\\hline
\verb|FindIf| & Copies the first matching value to \verb|result|.  Returns one
if found, zero if absent, or a negative error. \\\hline
\verb|AnyOf| & Returns one if any value matches, zero otherwise, or a negative
error. \\\hline
\verb|AllOf| & Returns one if every value matches, zero otherwise, or a
negative error.  It returns one for an empty range. \\\hline
\verb|CountIf| & Stores the number of matching values in \verb|*result| and
returns one.  The output is valid only after success. \\\hline
\verb|ToVector| & Copies every produced value into a new Vector allocated with
the range allocator.  The caller owns the Vector. \\\hline
\end{tabular}
\end{center}

\subsection{Pipeline example}

The following pipeline prints 4 and 16.  Each construction step is checked;
because adaptors are transactional, \verb|range| always remains valid for the
single cleanup at the end.

\begin{verbatim}
#include <stdio.h>
#include "range.h"

static int is_even(const void *element, void *arg)
{
    (void)arg;
    return (*(const int *)element % 2) == 0;
}

static int square(const void *input, void *output, void *arg)
{
    int value = *(const int *)input;
    (void)arg;
    *(int *)output = value * value;
    return 1;
}

static int print_int(const void *element, void *arg)
{
    (void)arg;
    printf("%d\n", *(const int *)element);
    return 1;
}

int main(void)
{
    int values[] = { 1, 2, 3, 4, 5, 6 };
    Range *range = NULL;
    int status;

    status = iRange.FromArray(values, 6, sizeof(int), &range);
    if (status > 0)
        status = iRange.Filter(&range, is_even, NULL);
    if (status > 0)
        status = iRange.Transform(&range, sizeof(int), square, NULL);
    if (status > 0)
        status = iRange.Take(&range, 2);
    if (status > 0)
        status = iRange.ForEach(range, print_int, NULL);

    if (range != NULL) {
        int finalizeStatus = iRange.Finalize(range);
        if (status >= 0)
            status = finalizeStatus;
    }
    return status < 0;
}
\end{verbatim}

\verb|GetElementSize| returns the current output element size, including the
new size established by \verb|Transform|.  \verb|Finalize| recursively
destroys the complete recipe.  Both operations reject invalid or stale inner
handles; this guards the unique-ownership rule described above.
